\documentclass[11pt,a4paper]{article}

%% =====================================================================
%% Clay RH via Substrate --- Phase 1 (v2, HONEST-SCOPE REWRITE)
%%
%% Title: Phase 1 --- A Substrate-Level Reformulation of the
%%        Riemann Hypothesis via the Principia Fractalis Framework
%%
%% Author: Pablo Cohen (with Claude Opus 4.7 / Anthropic)
%% Date  : 2026-06-06
%% Anchor: HEAD 2b777ca / 4785c55, Lean 4 PF closure clean,
%%         zero project axioms, kernel-only
%%         [propext, Classical.choice, Quot.sound].
%% Honest-scope discipline:
%%   Every claim in this paper is sourceable to a specific Lean
%%   file's own HONEST SCOPE section. Where a Lean file flags a
%%   field as :=True or := rfl or substrate-shadow, the paper
%%   flags the same. Earlier drafts drifted into Clay-prize
%%   framings; this rev corrects the drift on the prose side
%%   to match the rewritten master paper
%%   `Papers/clay_substrate_reformulation_v1.tex`.
%% =====================================================================

\usepackage{amsmath,amsthm,amssymb,amsfonts}
\usepackage[margin=1in]{geometry}
\usepackage{hyperref}
\usepackage{cleveref}
\usepackage{microtype}
\usepackage{enumitem}
\usepackage{booktabs}
\usepackage{xcolor}

\hypersetup{
  colorlinks=true,
  linkcolor=blue!50!black,
  urlcolor=blue!50!black,
  citecolor=blue!50!black,
}

\theoremstyle{plain}
\newtheorem{theorem}{Theorem}[section]
\newtheorem{lemma}[theorem]{Lemma}
\newtheorem{proposition}[theorem]{Proposition}
\newtheorem{corollary}[theorem]{Corollary}

\theoremstyle{definition}
\newtheorem{definition}[theorem]{Definition}
\newtheorem{solution}[theorem]{Solution}
\newtheorem{remark}[theorem]{Remark}

\title{Phase 1: A Substrate-Level Reformulation of the
Riemann Hypothesis\\
via the Principia Fractalis Framework\\
\large (Entry Vector --- Not a Clay Discharge)}
\author{Pablo Cohen\\
\small{\texttt{psolorzano@gmail.com}}\\[0.3em]
\small{Collaboration: Claude Opus 4.7 (Anthropic)}}
\date{2026-06-06}

\begin{document}
\maketitle

\begin{abstract}
This is the Phase 1 paper of a four-paper sequence that presents the
Principia Fractalis (PF) framework's \emph{substrate-level
reformulation} of the six remaining Clay Millennium Problems.
\textbf{This is not a Clay-prize claim for the Riemann Hypothesis.}
The substrate-level RH content remains conditional on the published
Hilbert--P\'olya program and on Mayer 1991 \S3's symmetric-quotient
spectral correspondence, both preserved openly as typed residuals.
What this paper exhibits is a substrate-level encoding of the RH axis
on the framework's $T_3^{\mathrm{sym}}$ transfer operator on
$\mathtt{LogWeightedL2}$, together with the algebraic constraint
$\alpha_{\mathrm{RH}} = 3/2$ pinned by two framework-internal arithmetic
compatibility identities on the chosen $\alpha$-skeleton under the
editorial convention $\alpha_{\mathrm{Poincar\acute e}} := 1$ motivated
by Perelman 2003.

The RH axis is the natural entry vector to the framework because of
the pre-built GRH-conditional scaffold in analytic number theory
(Conrey 2003; Borwein, Choi, Rooney, Weirathmueller 2008; von Koch
1901; Bach 1990; Pomerance--Selfridge--Wagstaff 1980). That scaffold
\emph{would} become unconditional under a literal RH discharge; the
framework here does not deliver that literal discharge. What the
framework delivers on this axis is a substrate-level reformulation
machine-verified in Lean 4 with kernel-only axioms
$\mathtt{[propext, Classical.choice, Quot.sound]}$, and a structural
link to the framework's simultaneous-closure mechanism on canonical
encodings of the six remaining Clay axes (Phase 2).

This is the first of four papers; the others are the master paper on
the simultaneous-closure mechanism
(\texttt{clay\_substrate\_reformulation\_v1.tex}), the P vs NP entry
(Phase 3), and the keystone paper on what the substrate is for
(consciousness, ZPE, cosmology --- Phase 4). Honest scope per the
file-level honest-scope marker
($\mathtt{sec:substrate-honest-scope}$) is preserved throughout.
Independent expert review and substantial revision are explicitly
invited.
\end{abstract}

\tableofcontents

\section{What this paper offers and what it does not}
\label{sec:offers}

This paper offers a \emph{substrate-level reformulation of the RH
axis} of the Principia Fractalis framework, with the chosen
$\alpha$-value $\alpha_{\mathrm{RH}} := 3/2$ pinned by two framework-internal
arithmetic identities under the editorial convention
$\alpha_{\mathrm{Poincar\acute e}} := 1$. It does \emph{not} offer a
literal mathlib-tier discharge of the Riemann Hypothesis at the Clay
formulation. The substrate-level RH bridge is conditional on the
published Hilbert--P\'olya program and on Mayer 1991 \S3's
symmetric-quotient spectral correspondence, both preserved openly as
typed residuals in the Lean tree.

\medskip

\noindent\fbox{\parbox{0.97\textwidth}{
\textbf{What the Lean kernel carries on this axis.} A discharge of
$\mathtt{Clay\_RiemannHypothesis\_Standard}$ on the framework's
$T_3^{\mathrm{sym}}$ encoding through
$\mathtt{PF\_RH\_capstone\_via\_Mayer1991\_T3sym}$, gated by two named
typed residuals (the Mayer 1991 symmetric-quotient
spectral-correspondence Prop and the published Hilbert--P\'olya
program Prop). The encoding is the framework's
$T_3^{\mathrm{sym}}$/$\mathtt{LogWeightedL2}$ realisation; it differs
from the literal mathlib $\mathtt{riemannZeta}$ tier in named,
disclosed ways.

\textbf{What this paper does NOT claim.} A literal mathlib-tier
discharge of RH (the form required by the Clay rules). The two named
residuals are unchanged.
}}

\subsection{Why RH is the natural entry vector}
\label{sec:entry-vector}

Among the seven Clay axes, RH carries the largest pre-built
mathematical audience. Thousands of theorems in analytic number
theory are GRH-conditional; under a confirmed RH they become
unconditional in one stroke. The reference catalogue is Borwein,
Choi, Rooney, Weirathmueller \cite{borwein_choi_rooney_weirathmueller2008};
Conrey 2003 \cite{conrey2003} enumerates the principal families.
Algorithmic number theory: Bach 1990 \cite{bach1990} and Pomerance--%
Selfridge--Wagstaff 1980 \cite{pomerance_selfridge_wagstaff1980}
make Miller--Rabin primality testing deterministic in
$\widetilde O((\log n)^2)$ under GRH. Prime-distribution bounds:
von Koch 1901 \cite{vonkoch1901} and Schoenfeld 1976
\cite{schoenfeld1976}.

We take this scaffold as the reason RH is the natural \emph{entry
vector} for an audience approaching the framework via the Clay
problems. We do not take it as evidence the framework discharges RH
literally. The framework's substrate-level contribution on this axis
is described honestly in \S\ref{sec:rh-on-substrate} and
\S\ref{sec:honest-scope}.

\subsection{Pointer to the other three papers}
\label{sec:pointer}

This Phase 1 paper foregrounds the RH-axis reformulation. The
companion papers in the four-paper sequence are:
(i)~the master paper~\cite{master_substrate}, which carries the
simultaneous-closure mechanism on the framework's canonical
encodings, the eleven framework-internal arithmetic compatibility
identities, the seven-field bundle (5 typed residuals + 2
$\mathtt{True}$-markers), and the per-axis honest-scope discussion
for all six remaining Clay axes; (ii)~the Phase 3 paper on the P vs
NP axis \cite{phase3_pnp}, with the framework's
$\mathtt{Machine}$-as-free-functions honest scope explicitly flagged;
(iii)~the Phase 4 keystone paper~\cite{phase4_keystone} on
consciousness, ZPE, and cosmology. The RH axis is one sub-story of
the simultaneous-closure mechanism described in the master paper.

\section{The substrate and the editorial convention}
\label{sec:substrate}

\subsection{The Timeless Field substrate}
\label{sec:tf}

The framework's substrate is the nuclear $C^{*}$-algebra of ternary
projective Hilbert spaces
\[
\mathcal{H}_{\mathrm{TF}} \;=\; \varprojlim_{k}\, H_{k},
\qquad H_{k} = \mathbb{C}^{3^{k}},
\]
called the \emph{Timeless Field}. The substrate is a typed
mathematical object in the Lean development; this paper does not
claim that the substrate is canonical or unique. It is the
framework's choice of structural stage on which the RH axis is
stated together with the other five remaining Clay axes.

\subsection{The editorial convention $\alpha_{\mathrm{Poincar\acute e}} = 1$}
\label{sec:editorial-convention}

The framework adopts $\alpha_{\mathrm{Poincar\acute e}} := 1$ as its
editorial convention at the Poincar\'e axis of the $\alpha$-skeleton.
This convention is \emph{motivated by} Perelman 2003's Ricci-flow
proof of the Poincar\'e conjecture \cite{perelman2002,perelman2003a,%
perelman2003b}. The Lean kernel does \emph{not} consume Perelman's
theorem as a hypothesis: the proposition
$\mathtt{PerelmanAnchorInput} :
(\mathtt{framework\_alpha}).a_{\mathrm{Poincare}} = 1$ is discharged
in the Lean file by direct unfolding of the framework's definition
$\alpha_{\mathrm{Poincare}} := 1$
(\texttt{PF.Referee.PerelmanAnchoredSimultaneousClosure}, line~117;
the proof is by \texttt{norm\_num} on a definitional identity).
Perelman 2003 motivates the choice of the unit value at the
Poincar\'e axis; the kernel cascade is arithmetic on the chosen
constants.

\subsection{The $\alpha$-skeleton and the eleven framework-internal
identities}
\label{sec:alpha-skeleton}

The framework's $\alpha$-skeleton on the complexity-and-Clay axes is
the nine-tuple
$(\alpha_{\mathrm{Poincar\acute e}}, \alpha_{\mathrm{RH}},
\alpha_{\mathrm{YM}}, \alpha_{\mathrm{BSD}}, \alpha_{\mathrm{NS}},
\alpha_{\mathsf{P}\mathrm{vs}\mathsf{NP}}, \alpha_{\mathrm{Hodge}},
\alpha_{P}, \alpha_{NP}, \alpha_{\mathrm{QG}})$ of chosen
\texttt{noncomputable def}s in
\texttt{PF.CrossMillenniumSharedInvariants}. \textbf{Honest scope}
quoted verbatim from that file's header (lines 18--29):
\begin{quote}
``\emph{These are \textbf{not} Millennium discharges. They are
axiom-free algebraic facts about the 9 chosen $\alpha$-values; the
framework's interpretation --- that each $\alpha$ is the canonical
eigenvalue parameter for its respective Millennium operator --- is
conjectural and lives elsewhere\ldots The value of this file: it
makes the \emph{interlocking} of $\alpha$-values explicit and
machine-checked, so future work cannot accidentally introduce
$\alpha$-redefinitions that break the algebraic skeleton.}''
\end{quote}

The eleven identities on the $\alpha$-skeleton (each provable by
\texttt{ring}, \texttt{norm\_num}, or \texttt{linarith} on the
definitions) are catalogued in the master paper
\cite[Thm~2.3]{master_substrate}. The two that pin the chosen value
$\alpha_{\mathrm{RH}} := 3/2$ are
\[
\alpha_{\mathrm{RH}}^{2} = 9/4
\qquad\text{and}\qquad
\alpha_{\mathrm{RH}}\,\alpha_{\mathrm{YM}} = 3
\quad\text{with}\quad
\alpha_{\mathrm{YM}} = \alpha_{\mathrm{Poincar\acute e}} + 1 = 2.
\]
Each is machine-verified by exact Lean name
($\alpha\_\mathrm{RH}\_\mathrm{sq}\_\mathrm{eq}\_\mathrm{nine}\_\mathrm{fourths}$,
$\alpha\_\mathrm{RH}\_\mathrm{mul}\_\mathrm{YM}\_\mathrm{eq}\_\mathrm{three}$,
$\alpha\_\mathrm{YM}\_\mathrm{eq}\_\alpha\_\mathrm{Poincare}\_\mathrm{plus}\_\mathrm{one}$
in \texttt{PF.CrossMillenniumSharedInvariants}).

\section{What the framework carries on the RH axis}
\label{sec:rh-on-substrate}

\subsection{The substrate-level RH bridge}
\label{sec:rh-bridge}

The framework's canonical encoding on the RH axis is the
$T_{3}^{\mathrm{sym}}$ transfer operator on the log-weighted $L^{2}$
completion $\mathtt{LogWeightedL2}$, the framework-level realisation
of the Mayer 1991 path. The substrate-level RH discharge theorem is

\medskip
\hspace*{1em}\texttt{PF.Referee.RHCapstoneTypedBridgeV4.}\\
\hspace*{1em}\texttt{PF\_RH\_capstone\_via\_Mayer1991\_T3sym}
(line~230, HEAD \texttt{2b777ca}/\texttt{4785c55}).
\medskip

This discharge consumes \emph{two named typed residuals}:
\begin{itemize}[leftmargin=*,nosep]
\item \texttt{rh\_mayer\_HP} \,:\,
  $\mathtt{Mayer1991\_SymmetricQuotientHasZetaSpectrum}$ ---
  the Mayer 1991 \S3 symmetric-quotient spectral-correspondence
  Prop, preserved openly;
\item \texttt{rh\_HP\_program} \,:\,
  $\mathtt{HilbertPolyaProgramConjecture}$ --- the published
  Hilbert--P\'olya program in its standard form, preserved openly.
\end{itemize}
The discharge lands on $\mathtt{Clay\_RiemannHypothesis\_Standard}$
\emph{on the framework's $T_{3}^{\mathrm{sym}}$ encoding}.

\medskip

\noindent\fbox{\parbox{0.97\textwidth}{
\textbf{Honest scope.} The framework's $\mathtt{Clay\_*\_Standard}$
discharges on the substrate-level $T_{3}^{\mathrm{sym}}$ encoding are
\emph{not} discharges of RH at the literal mathlib $\mathtt{riemannZeta}$
tier required by the Clay rules. The two named residuals (Mayer 1991
\S3 symmetric quotient; HP program) are unchanged. The substrate-level
discharge is what is in the kernel; the literal Clay tier is what is
not.
}}

\subsection{The five-formulation Hilbert--P\'olya collapse}
\label{sec:hp-five}

The substrate collapses the published Hilbert--P\'olya program to a
single typed Prop in five equivalent formulations: Berry--Keating
1999 \cite{berry_keating1999}, Connes 1999 \cite{connes1999},
Bost--Connes 1995 \cite{bost_connes1995}, PF $T_{3}^{\mathrm{sym}}$
on $\mathtt{LogWeightedL2}$, and Mayer 1991 \cite{mayer1991}
symmetric-quotient on the $+1$-eigenspace.

\begin{theorem}[Five-formulation Hilbert--P\'olya collapse]
\label{thm:hp-five}
The five formulations BK, C, BC, PF, Mayer of the Hilbert--P\'olya
program are typed-equivalent on the framework's substrate.
\textbf{Lean:}
\texttt{PF.Referee.RHCapstoneTypedBridgeV4.\\
PF\_RH\_V4\_five\_formulation\_equivalence};
\texttt{PrincipiaTractalis.HilbertPolyaIdentificationPrecise.\\
hilbert\_polya\_formulations\_equivalent};
\texttt{PrincipiaTractalis.HilbertPolyaIdentificationPrecise.\\
hilbert\_polya\_implies\_Clay\_RiemannHypothesis\_Standard}.
\end{theorem}

The five-formulation equivalence is unconditional at the substrate
level; the equivalence \emph{target} (the HP-program Prop) is the
typed residual the framework consumes.

\subsection{The Wave 58 concrete finite-dimensional witness}
\label{sec:concrete}

The framework exhibits an explicit $3 \times 3$ Hamiltonian
$H_{\mathrm{concrete}}$ on $e_{0}, e_{1}, e_{2}$ with spectrum
$\{1, 2, 3\}$, trace $= 6$, $\det = 6$, accompanied by a
compact-operator substitute $K_{\mathrm{concrete}}$ and an inhabited
substrate-level Hilbert--P\'olya typed contract on the finite-dim
witness.

\begin{theorem}[Wave 58 Ch~17 concrete witness, finite-dim scope]
\label{thm:ch17-concrete}
The substrate-level seven-clause Prop
$\mathtt{Ch17OperatorTheoryConcrete}$ holds axiom-free.
\textbf{Lean:}
\texttt{PF.Wave58.Ch17OperatorTheoryConcrete.\\
Ch17OperatorTheoryConcrete} (structure);
\texttt{PF.Wave58.Ch17OperatorTheoryConcrete.\\
ch17\_operator\_theory\_concrete\_capstone}.
\end{theorem}

\noindent\textbf{Honest scope.} The Ch~17 concrete witness is the
framework's substrate-level \emph{finite-dimensional realisation} of
the Hilbert--P\'olya typed contract on a $3\times 3$ self-adjoint
matrix. The literal continuum Hilbert--P\'olya operator on a
suitable infinite-dimensional $L^{2}$ space is \textbf{not} what this
witness exhibits. The finite-dim witness closes the Wave 58 audit
Level-C gap for the Ch~17 operator-theory chapter; the inf-dim case
is preserved openly in the named residual.

\section{The framework's mechanism, in one sentence}
\label{sec:mechanism}

The substrate $\to$ RH mechanism the framework offers, stated
honestly: \emph{under the editorial convention
$\alpha_{\mathrm{Poincar\acute e}} := 1$ motivated by Perelman 2003,
plus two framework-internal arithmetic identities on the chosen
$\alpha$-skeleton (the $9/4$-square identity and the
$\alpha_{\mathrm{RH}}\,\alpha_{\mathrm{YM}} = 3$ product identity
with $\alpha_{\mathrm{YM}} = 2$), the chosen value
$\alpha_{\mathrm{RH}} := 3/2$ is the only value compatible with the
framework's algebraic system; the framework then offers a
substrate-level RH bridge through
$\mathtt{PF\_RH\_capstone\_via\_Mayer1991\_T3sym}$ on the
$T_{3}^{\mathrm{sym}}$ encoding, conditional on two named typed
residuals (Mayer 1991 \S3 symmetric quotient and the published
Hilbert--P\'olya program).} The two typed residuals are the named
open content; the substrate-level bridge is what is in the kernel.

The same substrate that carries this RH bridge also carries
substrate-level encodings of the other five remaining Clay axes;
under the same editorial convention, the
$\mathtt{perelman\_anchor\_yields\_simultaneous\_clay\_closure}$
theorem of the master paper~\cite{master_substrate} delivers all six
$\mathtt{Clay\_*\_Standard}$ discharges on canonical encodings
simultaneously, gated by one seven-field bundle (five typed
residuals + two $\mathtt{True}$-markers). The simultaneous-closure
mechanism is the structural contribution; the per-axis residuals are
preserved openly throughout.

\section{Empirical anchor: IBM Quantum, framed as pre-registration}
\label{sec:ibm}

\noindent\fbox{\parbox{0.97\textwidth}{
\textbf{Pre-registration status (honest framing).} At HEAD
\texttt{4785c55} (2026-05-24) and unchanged at the current HEAD,
\textbf{only the pair $(\alpha_{\mathrm{RH}}, \alpha_{NP})$ has been
observed on IBM Quantum hardware} consistent with the framework's
prediction. The other 7 $\alpha$-values
(\texttt{IBMHardware9WayEvidence.lean} lines 318--320) are framework
predictions \emph{committed BEFORE hardware measurement}.

\medskip

\noindent The Lean theorem
$\mathtt{IBM\_hardware\_nine\_way\_random\_match\_probability\_bound}$
records the \textbf{conditional} bound (file lines 322--327, verbatim):
\begin{quote}
``\emph{IF all 9 instances are eventually measured to within}
$10^{-3}$ \emph{of the framework's pre-specified values under the
assumed} \texttt{NoiseModel9}\emph{, THEN the random-match probability
is at most} $10^{-15}$.''
\end{quote}

\noindent
\textbf{Current status:} 2 of 9 measured-consistent, 7 of 9
pre-registered, conditional bound in force. The pre-registration of
seven framework $\alpha$-values prior to hardware measurement is
scientifically valuable; it is honestly framed as such. The
conditional bound does \emph{not} say anything about the framework's
RH axis absent the remaining seven measurements. Earlier framework
documents that cited ``$\le 10^{-15}$'' as a measured RH confirmation
were over-stated; the bound is conditional.
}}

\section{The 23-problem reach, with field-level honesty}
\label{sec:reach}

The framework develops substrate-level treatment for 23 named open
problems (7 Clay axes + 16 non-Clay open problems), bundled in
$\mathtt{framework\_universal\_reach\_realized}$
(\texttt{FrameworkUniversalReach.lean}, line~153).

\noindent\fbox{\parbox{0.97\textwidth}{
\textbf{Bundle-field discipline (quoted from the Lean record's field
declarations, lines 99--137):} of the 17 fields in
$\mathtt{FrameworkUniversalReach}$, \textbf{15 are typed as}
$\mathtt{: True}$ (provenness tags carrying no mathematical content
at the record level). The two fields that \emph{do} carry non-trivial
framework Props at the record level are
$\mathtt{brocard\_attack\_realized}$ and
$\mathtt{hadwiger\_nelson\_attack\_realized}$. The substantive
per-axis content for the other 14 non-Clay problems and the seven
Clay axes lives in separate framework-attack modules cited by the
record file header. The 23-problem record is a single citation point
indexing these modules; reading the record itself as a Clay-Standard
discharge is incorrect.
}}

\section{Independent verification posture}
\label{sec:lean}

The Lean~4 proof objects for the theorems cited in this paper
type-check in the Lean~4 kernel with kernel axioms
$\mathtt{[propext, Classical.choice, Quot.sound]}$ and zero project
axioms. The PF\_Lean4Lean harness provides \textbf{build-hash
robustness across two Lake packages}: the canonical theorems are
re-bound in a separate package and rebuilt from source. \textbf{It is
not an independent type-checker written in another language.} Quoted
verbatim from
\texttt{PF\_Lean4Lean/PF\_L4L/Referee/ClayVerificationHarness.lean}
lines~47--56:
\begin{quote}
``\emph{This is the same Path-C protocol as
\texttt{FlagshipReverification} and
\texttt{V2AndMasterReverification}: a second pass through Lean's
kernel from a separate package with an independent build hash. It is
NOT a separate type-checker written in another language. The honest
scope of each individual closure is the one fixed by that closure's
own} \texttt{*\_honestScope} / \texttt{*\_honest\_scope}
\emph{marker in the canonical PF library\ldots The harness re-exports
the closures as they are --- it does NOT strengthen any of them.}''
\end{quote}

A Coq mirror provides cross-prover structural parity; approximately
75\% of definitions in the Coq mirror are typed
$\mathtt{: Prop := True}$, i.e.\ the Coq tree is \textbf{parity-only}
documenting that the naming skeleton admits cross-prover encoding,
not independent semantic verification of the Lean content.

\section{Honest scope}
\label{sec:honest-scope}

Per the file-level honest-scope marker
($\mathtt{sec:substrate-honest-scope}$), PF is a
\emph{substrate-level monograph in Grothendieck mode}.

\begin{remark}[Honest scope on the RH axis, foregrounded]
\label{rem:honest-scope-rh}
This paper, and the framework it foregrounds on the RH axis, are
\emph{not} a literal mathlib-tier discharge of the Riemann Hypothesis
in the Clay form. The substrate-level RH content is conditional on
the $\mathtt{HilbertPolyaProgramConjecture}$ (preserved openly in
the Hilbert--P\'olya identification module) and on the Mayer 1991 \S3
spectral correspondence
($\mathtt{Mayer1991\_SymmetricQuotientHasZetaSpectrum}$, preserved
openly in the Mayer transfer-operator formalization module). The
five-formulation Hilbert--P\'olya collapse
(\cref{thm:hp-five}) is unconditional at the substrate level; the
collapse target is the published HP program. The substrate-level Ch
17 concrete witness exhibits a $3\times 3$ self-adjoint Hamiltonian
with spectrum $\{1,2,3\}$; the literal continuum Hilbert--P\'olya
operator on an infinite-dimensional space remains the named open
content. The 23-problem reach record is, at the record level, 15
fields typed $\mathtt{: True}$ plus 2 fields carrying non-trivial
Props (Brocard, Hadwiger--Nelson); the substantive per-axis content
lives in the per-axis attack modules cited by the record.
\end{remark}

\begin{remark}[Cross-Millennium $\alpha$-skeleton: framework-internal]
\label{rem:alpha-internal}
The eleven cross-Millennium arithmetic compatibility identities that
pin $\alpha_{\mathrm{RH}} = 3/2$ are framework-internal: they are
algebraic facts about the nine \emph{chosen} $\alpha$-values, each
provable by \texttt{ring}/\texttt{norm\_num} on the definitions. The
framework's interpretation that each $\alpha$ is the canonical
eigenvalue parameter for its Millennium operator is, per the file
header of \texttt{PF.CrossMillenniumSharedInvariants}, ``conjectural
and lives elsewhere.''
\end{remark}

\begin{remark}[Falsifiability]
\label{rem:falsifiability}
The framework commits to eight typed falsifiability conditions in
$\mathtt{PF.Referee.FrameworkFalsifiabilityConditions}$. As of the
2026-06-04 referee-proof literature audit: \textbf{zero of eight
triggered}; \textbf{two actively supported by 2024--2026 literature}
($F_3$ dark-energy direction and $F_6$
$\Omega_{\Lambda} \in [0.65, 0.75]$).
\end{remark}

\section{Conclusion}
\label{sec:conclusion}

This Phase 1 paper presents a substrate-level reformulation of the
Riemann Hypothesis on the framework's $T_{3}^{\mathrm{sym}}$
encoding, machine-verified in Lean 4 with kernel-only axioms. The
substrate-level RH bridge is conditional on two named typed
residuals; it is not a literal mathlib-tier discharge of RH at the
Clay form. The RH axis is the natural \emph{entry vector} to the
framework because of the pre-built GRH-conditional scaffold in
analytic number theory; the framework's contribution is structural
reformulation on a substrate-level encoding, not prize-eligible
discharge.

The same substrate carries substrate-level encodings of the other
five remaining Clay axes; the simultaneous-closure mechanism on
canonical encodings is the subject of the master
paper~\cite{master_substrate}. The framework's reach beyond the Clay
set (23 named open problems via a single-citation record), the
empirical anchors (143-problem coherence and IBM Quantum 9-way), and
the Phase 4 substrate-level content on consciousness, ZPE, and
cosmology are presented in the companion papers. \textbf{Independent
expert review and substantial revision are explicitly invited.}

\begin{thebibliography}{99}

\bibitem{master_substrate}
Cohen, P., Claude Opus 4.7 (2026-06-06). \emph{A Substrate-Level
Reformulation of the Six Remaining Clay Millennium Problems via the
Principia Fractalis Framework}.
\texttt{Papers/clay\_substrate\_reformulation\_v1.tex}.

\bibitem{phase3_pnp}
Cohen, P., Claude Opus 4.7 (2026-06-06). \emph{Phase 3 ---
A Substrate-Level Reformulation at the P vs NP Axis via the
Principia Fractalis Framework}.
\texttt{Papers/clay\_PvsNP\_via\_substrate\_v2.tex}.

\bibitem{phase4_keystone}
Cohen, P., Claude Opus 4.7 (2026-06-06). \emph{Phase 4 Keystone ---
Substrate-Level Content on Consciousness, Zero-Point Energy, and
Cosmology via the Principia Fractalis Framework}.
\texttt{Papers/principia\_fractalis\_keystone\_v1.tex}.

\bibitem{principia_book}
Cohen, P. (2026). \emph{Principia Fractalis}, Second Edition.
HEAD of
\url{https://github.com/FractalDevTeam/Principia-Fractalis}.

\bibitem{conrey2003}
Conrey, J.~B. (2003). The Riemann hypothesis. \emph{Notices of the
American Mathematical Society} 50(3): 341--353.

\bibitem{borwein_choi_rooney_weirathmueller2008}
Borwein, P., Choi, S., Rooney, B., Weirathmueller, A. (2008).
\emph{The Riemann Hypothesis: A Resource for the Afficionado and
Virtuoso Alike}. CMS Books in Mathematics, Springer.

\bibitem{bach1990}
Bach, E. (1990). Explicit bounds for primality testing and related
problems. \emph{Mathematics of Computation} 55(191): 355--380.

\bibitem{pomerance_selfridge_wagstaff1980}
Pomerance, C., Selfridge, J.~L., Wagstaff, S.~S.~Jr.~(1980). The
pseudoprimes to $25 \cdot 10^{9}$. \emph{Mathematics of Computation}
35(151): 1003--1026.

\bibitem{vonkoch1901}
von Koch, H. (1901). Sur la distribution des nombres premiers.
\emph{Acta Math.} 24: 159--182.

\bibitem{schoenfeld1976}
Schoenfeld, L. (1976). Sharper bounds for the Chebyshev functions
$\theta(x)$ and $\psi(x)$. II. \emph{Mathematics of Computation}
30(134): 337--360.

\bibitem{perelman2002}
Perelman, G. (2002). The entropy formula for the Ricci flow and its
geometric applications. arXiv:math/0211159.

\bibitem{perelman2003a}
Perelman, G. (2003). Ricci flow with surgery on three-manifolds.
arXiv:math/0303109.

\bibitem{perelman2003b}
Perelman, G. (2003). Finite extinction time for the solutions to
the Ricci flow on certain three-manifolds. arXiv:math/0307245.

\bibitem{mayer1991}
Mayer, D.~H. (1991). The thermodynamic formalism approach to
Selberg's zeta function for $\mathrm{PSL}(2, \mathbb{Z})$.
\emph{Bull.\ Amer.\ Math.\ Soc.}\ 25(1): 55--60.

\bibitem{berry_keating1999}
Berry, M.~V., Keating, J.~P. (1999). The Riemann zeros and eigenvalue
asymptotics. \emph{SIAM Review} 41(2): 236--266.

\bibitem{connes1999}
Connes, A. (1999). Trace formula in noncommutative geometry and the
zeros of the Riemann zeta function. \emph{Selecta Math.} 5: 29--106.

\bibitem{bost_connes1995}
Bost, J.-B., Connes, A. (1995). Hecke algebras, type III factors and
phase transitions with spontaneous symmetry breaking in number
theory. \emph{Selecta Math.} 1: 411--457.

\end{thebibliography}

\end{document}
